\documentclass[12pt,a4paper]{article}
\usepackage[utf8]{inputenc}
\usepackage{amsmath,amssymb,amsthm}
\usepackage{geometry}
\usepackage{algorithm}
\usepackage{algpseudocode}
\geometry{margin=1in}

\newtheorem{theorem}{Theorem}
\newtheorem{lemma}[theorem]{Lemma}
\newtheorem{proposition}[theorem]{Proposition}
\newtheorem{corollary}[theorem]{Corollary}
\theoremstyle{definition}
\newtheorem{definition}[theorem]{Definition}
\theoremstyle{remark}
\newtheorem*{remark}{Remark}

\title{Problem 60: Prime Pair Sets}
\author{Project Euler Solution}
\date{}

\begin{document}
\maketitle

\section{Problem Statement}
Find the lowest sum for a set of five primes for which any two primes concatenate (in either order) to produce another prime.

\section{Mathematical Development}

\begin{definition}[Concatenation]
For positive integers $a$ and $b$, define
\[
a \| b = a \cdot 10^{d(b)} + b,
\]
where $d(b) = \lfloor \log_{10} b \rfloor + 1$ is the digit count of $b$.
\end{definition}

\begin{definition}[Prime Pair]
Primes $p$ and $q$ form a \emph{prime pair}, written $p \sim q$, if both $p \| q$ and $q \| p$ are prime.
\end{definition}

\begin{definition}[Prime Pair Graph]
The \emph{prime pair graph} is $G = (V, E)$ where $V$ is a set of primes and $\{p, q\} \in E$ iff $p \sim q$.
\end{definition}

\begin{theorem}[Modular Necessary Condition]
\label{thm:mod3}
Let $p, q > 3$ be distinct primes with $p \not\equiv q \pmod{3}$. Then $p \not\sim q$.
\end{theorem}

\begin{proof}
Since $10 \equiv 1 \pmod{3}$:
\[
p \| q = p \cdot 10^{d(q)} + q \equiv p + q \pmod{3}.
\]
If $p \equiv 1$ and $q \equiv 2 \pmod{3}$ (or vice versa), then $p + q \equiv 0 \pmod{3}$, making $p \| q$ divisible by~3. Since $p \| q > 3$, it is composite. The same argument applies to $q \| p$.
\end{proof}

\begin{corollary}[Residue Class Constraint]
In any prime pair set of primes $> 3$, all elements share the same residue modulo~3.
\end{corollary}

\begin{lemma}[Special Role of 3]
The prime 3 can pair with primes of either residue class modulo~3, since $3 + p \equiv p \pmod{3} \not\equiv 0$ for $p > 3$.
\end{lemma}

\begin{theorem}[Graph-Theoretic Reduction]
Finding a prime pair set of size~5 with minimum sum is equivalent to finding a minimum-weight 5-clique in the prime pair graph.
\end{theorem}

\begin{proof}
A prime pair set of size~5 is, by definition, a set of 5 pairwise-adjacent vertices in $G$, i.e., a 5-clique. The weight is the sum of the vertex labels.
\end{proof}

\begin{lemma}[Search Bound]
It suffices to search primes below $B = 10{,}000$. If a clique with sum $\Sigma$ is found, any clique containing a prime $> \Sigma$ has sum $> \Sigma$.
\end{lemma}

\begin{lemma}[Primality of Concatenations]
For primes below $10{,}000$, all concatenations are below $10^8$. Deterministic Miller--Rabin with witnesses $\{2, 3, 5, 7, 11, 13\}$ is correct for all $n < 3.2 \times 10^{14}$.
\end{lemma}

\begin{theorem}[Incremental Clique Search]
The 5-clique search proceeds by depth-first enumeration in increasing order with:
\begin{enumerate}
    \item Adjacency-intersection pruning at each depth.
    \item Branch-and-bound: prune when partial sum plus minimum remaining contribution exceeds best known.
    \item Mod-3 filtering per Theorem~\ref{thm:mod3}.
\end{enumerate}
\end{theorem}

\section{Algorithm}

\begin{algorithm}[H]
\caption{Find Minimum-Weight 5-Clique}
\begin{algorithmic}[1]
\State Sieve primes below $B$; build adjacency lists via pairwise primality tests
\State $\textit{best} \gets \infty$
\For{each $p_1$}
    \State $S_1 \gets \{p > p_1 : p_1 \sim p\}$
    \For{each $p_2 \in S_1$}
        \State $S_2 \gets \{p \in S_1 : p > p_2,\; p_2 \sim p\}$
        \For{each $p_3 \in S_2$}
            \State $S_3 \gets \{p \in S_2 : p > p_3,\; p_3 \sim p\}$
            \For{each $p_4 \in S_3$}
                \State $S_4 \gets \{p \in S_3 : p > p_4,\; p_4 \sim p\}$
                \For{each $p_5 \in S_4$}
                    \State $\textit{best} \gets \min(\textit{best},\; p_1 + p_2 + p_3 + p_4 + p_5)$
                \EndFor
            \EndFor
        \EndFor
    \EndFor
\EndFor
\State \Return $\textit{best}$
\end{algorithmic}
\end{algorithm}

\section{Complexity Analysis}

\begin{theorem}[Pair Precomputation]
Building the adjacency structure requires $O(P^2 \log^2 B)$ time, where $P = \pi(B) = 1229$.
\end{theorem}

\begin{proof}
There are $\binom{P}{2} = O(P^2)$ pairs. Each pair test performs two concatenations and two Miller--Rabin tests, each costing $O(\log^2 B)$.
\end{proof}

\begin{theorem}[Clique Search]
The worst-case clique search is $O(P^5)$, but adjacency-intersection and branch-and-bound pruning reduce the practical cost far below this bound. The dominant cost is the $O(P^2)$ pair precomputation.
\end{theorem}

\noindent\textbf{Space:} $O(P^2)$ for the adjacency structure.

\section{Result}
The answer is $\boxed{26033}$, from the set $\{13, 5197, 5701, 6733, 8389\}$.

\section{Answer}

\[
\boxed{26033}
\]

\end{document}
